% rel_chapter_4_sec45-46.tex
% Agent 18: Relativistic Q-law and Overstability (Ch IV §45-46)
% BDNK conversion: IS relaxation replaced by BDNK first-order
%   constitutive relations (Bemfica et al. 2018, Kovtun 2019)
% Conventions: (-,+,+,+) signature, c kept explicit, see RELATIVISTIC_CONVENTIONS.md

\section{The $\pi^{2}Q$-law and invariant --- relativistic derivation}
\label{sec:rel-45}

In the classical theory (\S\,45) Chandrasekhar showed that for large
Chandrasekhar numbers $Q\to\infty$, the critical Rayleigh number for the
onset of stationary convection obeys $R_{c}\to\pi^{2}Q$, independently of
the boundary conditions.  The physical origin of this law is the dominance
of Ohmic (Joule) dissipation over viscous dissipation in the limit of
strong magnetic fields or vanishing viscosity.  We now re-derive these
results within the relativistic magneto-hydrodynamic framework, where the
role of inertia is played by the \emph{relativistic enthalpy density}
$w = \varepsilon + p$ rather than the rest-mass density $\rho_{0}$ alone,
and where the magnetic field enters through the covariant description in
terms of the magnetic four-vector $b^{\mu}$.

%------------------------------------------------------------------
\subsection{Relativistic MHD perturbation equations for the inviscid limit}
\label{subsec:rel-45-inviscid}

In the relativistic framework the equilibrium is characterised by the
enthalpy density $w_{0} = \varepsilon_{0} + p_{0}$, the background
magnetic four-vector $b^{\mu}_{0}$ with $b^{2}_{0} = b^{\mu}_{0}\,b_{0\mu}$,
the thermal diffusivity $\kappa$, and the magnetic diffusivity
$\eta = c^{2}/(4\pi\sigma_{\mathrm{e}})$, where $\sigma_{\mathrm{e}}$ is
the electrical conductivity.  The relativistic Chandrasekhar number is
\begin{equation}
  \Qrel = \frac{b^{2}_{0}\,d^{2}}{4\pi\,(w_{0}/c^{2})\,\nu\,\eta}
        = \frac{\mu H^{2} d^{2}}{\rho_{0}\,\nu\,\eta}
          \left[\frac{\rho_{0}\,c^{2}}{w_{0}}\right] ,
  \label{eq:rel-Qrel-def}
\end{equation}
where the factor $\rho_{0}c^{2}/w_{0} = 1/h$ (with $h$ the specific
enthalpy) encodes the relativistic correction.  In the non-relativistic
limit $h\to 1$ and $\Qrel\to Q$.

For a fluid layer of depth $d$ with an imposed vertical magnetic field,
the linearised perturbation equations for the vertical velocity $W$,
the temperature perturbation $\Theta$, and the magnetic scalar $K$
follow from the relativistic MHD energy-momentum conservation
$\nabla_{\nu}T^{\mu\nu}_{\mathrm{total}} = 0$ together with the
induction equation.  In the inviscid limit ($\nu\to 0$) and for
stationary onset ($p = 0$), the dissipative terms vanish
and the perturbation equations reduce to the \emph{same structure}
as in the classical case, up to the replacement of $\rho_{0}$ by the
effective inertial mass density $w_{0}/c^{2}$.  This identity at
$\sigma=0$ holds for any causal dissipation formalism---whether
Israel--Stewart or BDNK---since the dissipative contributions are
proportional to~$\sigma$.  Specifically, in
dimensionless form (with $d$ as the unit of length) the vertical
velocity satisfies
\begin{equation}
  D^{2}(D^{2}-a^{2})W = \frac{\Rarel}{\Qrel}\,a^{2}\,W ,
  \label{eq:rel-inviscid-W}
\end{equation}
with boundary conditions $W = 0$ and $D^{2}W = 0$ at $z = 0$ and $z=1$.
Here $D = \dd/\dd z$, $a$ is the horizontal wave number, and
\begin{equation}
  \Rarel = \frac{g\,\alpha\,\beta\,d^{4}}{\nu\,\kappa}
           \left[\frac{\rho_{0}\,c^{2}}{w_{0}}\right]
  \label{eq:rel-Rarel-Q}
\end{equation}
is the relativistic Rayleigh number (cf.\ RELATIVISTIC\_CONVENTIONS.md).

The structural identity of~\eqref{eq:rel-inviscid-W} with the classical
equation~(4-197) is a direct consequence of the fact that in a static
equilibrium, the relativistic effects modify only the definition of
the dimensionless groups and not the operator structure.

%------------------------------------------------------------------
\subsection{Origin of the law: magnetic tension vs.\ relativistic inertia}
\label{subsec:rel-45-origin}

The lowest eigenfunction of~\eqref{eq:rel-inviscid-W} is again
$W = A\sin\pi z$, yielding the characteristic equation
\begin{equation}
  \frac{\Rarel}{\Qrel}
  = \frac{\pi^{2}(\pi^{2}+a^{2})}{a^{2}} .
  \label{eq:rel-RoverQ}
\end{equation}
The minimum occurs as $a\to\infty$:
\begin{equation}
  \boxed{\;\Rarel \;\to\; \pi^{2}\,\Qrel \quad (a\to\infty).}
  \label{eq:rel-pi2Q-law}
\end{equation}

The physical interpretation is entirely analogous to the classical one,
but cast in relativistic language:
\begin{itemize}
  \item The \emph{restoring force} is the magnetic tension
    $\sim b^{2}_{0}/(4\pi)$, which opposes bending of field lines.
  \item The \emph{inertia} resisting convective motions is
    $w_{0}/c^{2} = (\varepsilon_{0}+p_{0})/c^{2}$, not merely
    $\rho_{0}$.  In a relativistically hot plasma (radiation-dominated,
    or with strong internal energy), $w_{0} \gg \rho_{0}\,c^{2}$ and
    the effective inertia is enhanced, meaning that a stronger
    temperature gradient is needed to overcome the magnetic suppression.
  \item The critical temperature gradient becomes
    \begin{equation}
      g\,\alpha\,\beta_{c}
      = \pi^{2}\,
        \frac{b_{0}^{2}\,\sigma_{\mathrm{e}}}{w_{0}/c^{2}}\,
        \kappa\,d^{-4}
      = \pi^{2}\,
        \frac{\mu H^{2}\sigma_{\mathrm{e}}}{4\pi\,\rho_{0}\,h}\,
        \kappa\,d^{-4} ,
      \label{eq:rel-beta-c-inviscid}
    \end{equation}
    which differs from the classical result~(4-202) by the factor
    $1/h$.  This is a genuine relativistic effect: the enhanced inertia
    makes the fluid \emph{harder} to destabilise.
\end{itemize}

In complete analogy with the classical discussion, including the viscous
correction at one higher order yields
\begin{equation}
  \Rarel = \frac{\pi^{2}+a^{2}}{a^{2}}\,
           \bigl(a^{4} + \Qrel\,\pi^{2}\bigr) ,
  \label{eq:rel-higher-order}
\end{equation}
and the critical wave number at marginal stability scales as
\begin{equation}
  a_{c}^{(\mathrm{rel})} \to
  \bigl(\tfrac{1}{2}\,\pi^{4}\,\Qrel\bigr)^{1/6}
  \quad (\Qrel\to\infty) .
  \label{eq:rel-ac-asymptotic}
\end{equation}
The effective kinematic viscosity imparted by the magnetic field is
\begin{equation}
  \nu_{\mathrm{eff}}^{(\mathrm{rel})}
  = \mathrm{const.}\;\frac{(\mu H d)^{2}}{w_{0}/c^{2}}
  = \frac{\rho_{0}\,c^{2}}{w_{0}}\;\nu_{\mathrm{eff}}^{(\mathrm{class})} ,
  \label{eq:rel-nu-eff}
\end{equation}
which is \emph{reduced} relative to the classical value by the factor
$1/h$, reflecting the greater difficulty of accelerating a relativistic
fluid.

%------------------------------------------------------------------
\subsection{The invariant in the relativistic eigenvalue problem}
\label{subsec:rel-45-invariant}

In the classical theory, Chandrasekhar identified the invariant
(eq.~(4-215)--(4-219)):
\[
  \Psi^{(\mathrm{class})}
  = \eta\times\text{(magnetic energy in }\mathbf{h}_{1}^{(0)})
  + \nu\times\text{(kinetic energy in }\mathbf{u}_{\perp})
  = \tfrac{1}{2}\rho_{0}\,\nu\,W_{0}^{2} ,
\]
independent of $Q$.

In the relativistic formulation, the energy densities must be
computed from the full energy-momentum tensor.  The magnetic energy
density in the fluid frame is $b_{1}^{2}/(8\pi)$ where $b_{1}^{\mu}$
is the perturbed magnetic four-vector, and the kinetic energy density
of horizontal motions is $\tfrac{1}{2}(w_{0}/c^{2})\,|\mathbf{u}_{\perp}|^{2}$.
The relativistic invariant is therefore
\begin{equation}
  \Psi^{(\mathrm{rel})}
  = \eta\,\frac{1}{8\pi}
    \langle\langle |b_{1\perp}|^{2}\rangle\rangle
  + \nu\,\frac{w_{0}}{2\,c^{2}}
    \langle\langle |\mathbf{u}_{\perp}|^{2}\rangle\rangle ,
  \label{eq:rel-invariant-def}
\end{equation}
where $\langle\langle\cdot\rangle\rangle$ denotes averaging over both
the horizontal plane and the vertical direction.

Using the eigenfunctions for two free boundaries and the relation between
$a$ and $\Qrel$ at marginal stability, one finds by a computation
entirely parallel to the classical one that
\begin{equation}
  \boxed{\;
  \Psi^{(\mathrm{rel})}
  = \frac{1}{2}\,\frac{w_{0}}{c^{2}}\,\nu\,W_{0}^{2}
  = \frac{1}{2}\,\rho_{0}\,h\,\nu\,W_{0}^{2} ,\;}
  \label{eq:rel-invariant-result}
\end{equation}
independently of $\Qrel$.  This reduces to the classical invariant
$\tfrac{1}{2}\rho_{0}\nu W_{0}^{2}$ when $h\to 1$.

\begin{relcorrection}
The relativistic invariant~\eqref{eq:rel-invariant-result} differs from
its classical counterpart by the factor $h = w_{0}/(\rho_{0}c^{2})$.
Physically, the invariance reflects a balance between magnetic energy
(which is frame-independent to leading order) and kinetic energy
(which involves the full relativistic inertia $w_{0}/c^{2}$).  In a
radiation-dominated plasma where $w_{0}\approx 4p_{0}$, the invariant
can be several times larger than the classical prediction at the same
$W_{0}$.
\end{relcorrection}


%==================================================================
\section{Overstability (oscillatory onset) with magnetic field ---
         relativistic treatment}
\label{sec:rel-46}

We now turn to the question of whether instability can arise as
oscillations of increasing amplitude (overstability) in the relativistic
regime.  The classical treatment (\S\,46) showed that for two free
boundaries and $p_{2} > p_{1}$ (i.e.\ $\kappa > \eta$), oscillatory
onset is preferred at sufficiently large~$Q$.  The relativistic
generalisation within the BDNK formalism introduces two features:
\begin{enumerate}
  \item The BDNK first-order constitutive relations
    ($\Pi^{\mu\nu}=-2\shearvisc\,\sigma^{\mu\nu}$,
    $Q^{\mu}=-\thermcond T\Delta^{\mu\alpha}\partial_{\alpha}(\mu/T)
    +\text{frame terms}$)
    do \emph{not} raise the order of the characteristic equation beyond
    the classical (Navier--Stokes) form.  The dispersion relation has the
    same number of roots as in the classical theory---no spurious
    relaxation modes appear.
  \item \emph{Causality} imposes an upper bound on the phase velocity
    of any oscillatory mode: $v_{\mathrm{ph}} \le c$, guaranteed by
    the BDNK coupled inequalities on transport and frame coefficients
    \citep{Bemfica2018,Kovtun2019}.
\end{enumerate}

%------------------------------------------------------------------
\subsection{Relativistic perturbation equations with BDNK
            dissipation}
\label{subsec:rel-46-perturb}

The linearised equations for the perturbations
$(W, \Theta, K)$
about a static equilibrium with a vertical magnetic field are obtained
from $\covdu T^{\mu\nu}_{\mathrm{total}} = 0$ supplemented by the
BDNK first-order constitutive relations
\citep{Bemfica2018,Kovtun2019,HoultKovtun2020}:
\begin{align}
  Q^{\mu}
  &= -\thermcond\,T\,\Delta^{\mu\nu}\partial_{\nu}(\mu/T)
     + \text{frame corrections} ,
  \label{eq:rel-BDNK-heat}\\[4pt]
  \Pi^{\mu\nu}
  &= -2\,\shearvisc\,\sigma^{\mu\nu} ,
  \label{eq:rel-BDNK-shear}\\[4pt]
  \Pi_{\mathrm{bulk}}
  &= -\bulkvisc\,\nabla_{\mu}u^{\mu} .
  \label{eq:rel-BDNK-bulk}
\end{align}
These are \emph{algebraic} relations (no time derivatives of the
dissipative fluxes), so they do not introduce additional dynamical
degrees of freedom.  When linearised and decomposed into normal modes
$\sim \exp(\sigma t + i\,\mathbf{k}\cdot\mathbf{x})$, the dissipative
terms enter at the same order as in the classical Navier--Stokes theory.
Causality is guaranteed not by relaxation times but by the BDNK frame
coefficients, which ensure strong hyperbolicity of the full system.

%------------------------------------------------------------------
\subsection{The relativistic characteristic equation}
\label{subsec:rel-46-chareq}

For two free boundaries with $W = A\sin\pi z$, the elimination of
$\Theta$, $K$, and the dissipative perturbations from the linearised
system yields the characteristic equation.  Let
\begin{equation}
  x = \frac{a^{2}}{\pi^{2}},\quad
  i\sigma_{1} = \frac{\sigma}{\pi^{2}},\quad
  R_{1} = \frac{\Rarel}{\pi^{4}},\quad
  Q_{1} = \frac{\Qrel}{\pi^{2}},
  \label{eq:rel-scaled-vars}
\end{equation}
The BDNK characteristic equation takes the form
\begin{multline}
  (1+x)\bigl(1+x+i\,h\,\sigma_{1}\bigr)
  \Bigl[\bigl(1+x+i\,p_{2}^{(\mathrm{rel})}\sigma_{1}\bigr)
        \bigl(1+x+i\,\sigma_{1}\bigr)
  + Q_{1}\Bigr] \\
  = R_{1}\,x\,
    \bigl(1+x+i\,p_{2}^{(\mathrm{rel})}\sigma_{1}\bigr) ,
  \label{eq:rel-char-full}
\end{multline}
where $h = w_{0}/(\rho_{0}c^{2})$ is the specific enthalpy and
\begin{equation}
  p_{1}^{(\mathrm{rel})} = \frac{\kappa}{\nu\,h} ,\quad
  p_{2}^{(\mathrm{rel})} = \frac{\eta}{\nu\,h} ,
  \label{eq:rel-p1p2}
\end{equation}
are the relativistic Prandtl ratios.

\begin{relcorrection}
Equation~\eqref{eq:rel-char-full} has the \emph{same order} in
$\sigma_{1}$ as its classical counterpart~(4-222): it is fourth-order
in~$\sigma_{1}$.  This is a key advantage of the BDNK formalism over
Israel--Stewart: no spurious relaxation modes appear, and the dispersion
relation contains only the physical hydrodynamic modes.  Causality is
ensured not by additional dynamical degrees of freedom but by the BDNK
coupled inequalities on transport coefficients and frame coefficients
\citep{Bemfica2018,Kovtun2019,BemficaDisconziNoronhaKovtun2023}.
\end{relcorrection}

%------------------------------------------------------------------
\subsection{Relation to the classical characteristic equation}
\label{subsec:rel-46-classical}

Equation~\eqref{eq:rel-char-full} is structurally identical to the
classical equation~(4-224) with the replacements
$p_{1}\to p_{1}/h$ and $p_{2}\to p_{2}/h$ and an overall factor of $h$
multiplying the inertial term.  In the non-relativistic limit $h\to 1$
the classical equation~(4-224) is recovered exactly.  Unlike the
Israel--Stewart formulation, no additional ``Eckart limit'' needs to be
taken---the BDNK characteristic equation is already of the same
polynomial order as the classical one.

%------------------------------------------------------------------
\subsection{Critical conditions: overstability vs.\ stationary onset}
\label{subsec:rel-46-critical}

To determine when overstability is preferred, we seek solutions of the
characteristic equation with $\sigma_{1}$ real (purely oscillatory
marginal state).

\subsubsection{BDNK critical conditions}

Separating real and imaginary parts of~\eqref{eq:rel-char-full} with
$\sigma_{1}\in\mathbb{R}$, one obtains the analogue of the classical
equations~(4-226)--(4-231):
\begin{equation}
  p_{1}^{(0)\,2}\,\sigma_{1}^{2}
  = \frac{p_{2}^{(\mathrm{rel})}-p_{1}^{(\mathrm{rel})}}{1+p_{2}^{(\mathrm{rel})}}\,Q_{1}
    - (1+x)^{2} ,
  \label{eq:rel-sigma1-sq}
\end{equation}
and
\begin{equation}
  R_{1}^{(\mathrm{ov})}
  = \frac{(1+p_{2}^{(\mathrm{rel})})(p_{1}^{(\mathrm{rel})}+p_{2}^{(\mathrm{rel})})}{p_{1}^{(0)\,2}}
    \,\frac{1+x}{x}\,
    \left[(1+x)^{2}
    + Q_{1}\,\frac{p_{1}^{(0)\,2}}{(1+p_{1}^{(\mathrm{rel})})(p_{1}^{(\mathrm{rel})}+p_{2}^{(\mathrm{rel})})}\right] .
  \label{eq:rel-R1-overstab}
\end{equation}

The condition for overstability to exist is $\sigma_{1}^{2} > 0$, which
requires
\begin{equation}
  p_{2}^{(\mathrm{rel})} > p_{1}^{(\mathrm{rel})}
  \quad\Longleftrightarrow\quad
  \frac{\eta}{\nu h} > \frac{\kappa}{\nu h}
  \quad\Longleftrightarrow\quad
  \eta > \kappa .
  \label{eq:rel-overstab-condition}
\end{equation}
This is the \emph{same} condition as in the classical theory ($\kappa <
\eta$), since the factor $1/h$ cancels.  Physically, the ratio of thermal
to magnetic diffusivity governs whether oscillatory modes can be sustained,
and this ratio is unaffected by the relativistic enhancement of inertia.

\emph{However}, the critical $Q$ at which overstability becomes the
preferred mode, $Q^{(\mathrm{tr})}$, is shifted.  Since
$p_{1}^{(\mathrm{rel})} = \kappa/(\nu h)$ and $p_{2}^{(\mathrm{rel})} = \eta/(\nu h)$, the
transition Chandrasekhar number satisfies
\begin{equation}
  Q_{\mathrm{rel}}^{(\mathrm{tr})}
  = h^{2}\;Q_{\mathrm{class}}^{(\mathrm{tr})}
    \bigl[1 + \order{v^{2}/c^{2}}\bigr] ,
  \label{eq:rel-Qtr}
\end{equation}
so that in a relativistically hot plasma ($h>1$), a \emph{stronger}
magnetic field is needed before oscillatory onset is preferred.

\subsubsection{Role of the Prandtl and magnetic Prandtl numbers}

Define the relativistic Prandtl number and magnetic Prandtl number:
\begin{equation}
  \mathrm{Pr}_{\mathrm{rel}} = \frac{\nu\,h}{\kappa}
  = h\,\mathrm{Pr}_{\mathrm{class}} ,\qquad
  \mathrm{Pm}_{\mathrm{rel}} = \frac{\nu\,h}{\eta}
  = h\,\mathrm{Pm}_{\mathrm{class}} .
  \label{eq:rel-Pr-Pm}
\end{equation}
In terms of these, the condition for overstability becomes simply
\begin{equation}
  \mathrm{Pr}_{\mathrm{rel}} > \mathrm{Pm}_{\mathrm{rel}}
  \quad\Longleftrightarrow\quad
  \kappa < \eta ,
  \label{eq:rel-overstab-Pr-Pm}
\end{equation}
and the asymptotic ($\Qrel\to\infty$) behaviours are
\begin{align}
  R_{c}^{(\mathrm{ov,rel})}
  &\to \pi^{2}\,
       \frac{(1+\mathrm{Pm}_{\mathrm{rel}}^{-1})}
            {(1+\mathrm{Pr}_{\mathrm{rel}}^{-1})\,
             \mathrm{Pm}_{\mathrm{rel}}^{-1}}\;\Qrel ,
  \label{eq:rel-Rc-ov-asymp}\\[4pt]
  a_{c}^{(\mathrm{ov,rel})}
  &\to \left[\frac{\pi^{4}}{2}\,
       \frac{\mathrm{Pr}_{\mathrm{rel}}^{-2}}
            {(1+\mathrm{Pr}_{\mathrm{rel}}^{-1})
             (\mathrm{Pr}_{\mathrm{rel}}^{-1}+\mathrm{Pm}_{\mathrm{rel}}^{-1})}
       \;\Qrel\right]^{1/6} .
  \label{eq:rel-ac-ov-asymp}
\end{align}

The gyration frequency at marginal overstability for $\Qrel\to\infty$ is
\begin{equation}
  |p|^{(\mathrm{rel})} \to
  \left[\frac{\mathrm{Pm}_{\mathrm{rel}}^{-1}
              -\mathrm{Pr}_{\mathrm{rel}}^{-1}}
             {\mathrm{Pm}_{\mathrm{rel}}^{-1}
              (1+\mathrm{Pr}_{\mathrm{rel}}^{-1})}
  \right]^{1/2}
  \frac{\vA^{(\mathrm{rel})}}{d} ,
  \label{eq:rel-gyration}
\end{equation}
where the relativistic Alfv\'en speed is
\begin{equation}
  \vA^{(\mathrm{rel})}
  = \frac{b_{0}\,c}{\sqrt{4\pi\,w_{0} + b_{0}^{2}}}
  \;\le\; c .
  \label{eq:rel-Alfven}
\end{equation}
In the non-relativistic limit $w_{0}\to\rho_{0}c^{2}$ and
$b_{0}^{2}\ll 4\pi w_{0}$, so
$\vA^{(\mathrm{rel})}\to H\sqrt{\mu/(4\pi\rho_{0})} = V_{A}$, recovering
the classical result~(4-243).

%------------------------------------------------------------------
\subsection{Comparison: classical vs.\ relativistic critical curves}
\label{subsec:rel-46-comparison}

\begin{relcorrection}
The key differences between the classical and relativistic critical curves
in the $(R_{c}, Q)$-plane are:

\begin{enumerate}
  \item \textbf{Stationary branch:} Both follow $R_{c}\sim\pi^{2}Q$ for
    large~$Q$, but with different proportionality constants.  The
    relativistic curve lies above the classical one by a factor~$h$
    (the increased inertia raises the critical $R$).

  \item \textbf{Overstable branch:} The overstable branch is also shifted
    upward by a factor involving~$h$.  The branch point $Q^{(\mathrm{tr})}$
    at which overstability becomes preferred moves to
    $Q_{\mathrm{rel}}^{(\mathrm{tr})} \approx h^{2}\,
    Q_{\mathrm{class}}^{(\mathrm{tr})}$.

  \item \textbf{Wave numbers:} The critical wave numbers $a_{c}$ at onset
    retain the same power-law dependence $\sim Q^{1/6}$ but with $\Qrel$
    in place of $Q$, giving slightly different numerical prefactors.

  \item \textbf{Oscillation frequency:} The gyration frequency of
    overstable oscillations is controlled by the relativistic Alfv\'en
    speed $\vA^{(\mathrm{rel})} < V_{A}$, so that relativistic oscillation
    frequencies are \emph{lower} than their classical counterparts at the
    same field strength.
\end{enumerate}
\end{relcorrection}

%------------------------------------------------------------------
\subsection{BDNK causality constraints on overstable modes}
\label{subsec:rel-46-BDNK-causality}

In the BDNK formalism, the characteristic equation~\eqref{eq:rel-char-full}
has \emph{no additional roots} beyond the standard hydrodynamic modes.
Causality is instead enforced by coupled inequalities on the transport
coefficients $(\shearvisc, \bulkvisc, \thermcond)$ and the BDNK frame
coefficients $(\varepsilon_1, \beta_1, \alpha_1)$:

\begin{enumerate}
  \item \textbf{Same polynomial order as classical:} The characteristic
    equation remains fourth-order in $\sigma_{1}$, with no spurious
    relaxation modes.  Every root corresponds to a physical hydrodynamic
    mode.

  \item \textbf{Frequency-independent Prandtl ratios:} The relativistic
    Prandtl ratios $p_{1}^{(\mathrm{rel})} = \kappa/(\nu h)$ and
    $p_{2}^{(\mathrm{rel})} = \eta/(\nu h)$ are constants (not
    frequency-dependent), simplifying the overstability analysis.

  \item \textbf{Causality conditions:} The BDNK coupled inequalities
    \citep{BemficaDisconziNoronhaKovtun2023} require
    \begin{equation}
      \shearvisc > 0, \quad
      \bulkvisc + \tfrac{2}{3}\shearvisc > 0, \quad
      \text{frame coefficients satisfy specific bounds},
      \label{eq:rel-BDNK-causality-conditions}
    \end{equation}
    which replace the Israel--Stewart requirements $\taupi>0$, $\tauq>0$,
    $\tauPi>0$.  These ensure that all characteristic speeds of the
    full nonlinear system are bounded by~$c$.
\end{enumerate}

%------------------------------------------------------------------
\subsection{Causality constraint on oscillatory modes}
\label{subsec:rel-46-causality}

\begin{causalitycheck}
Every oscillatory mode at marginal stability has a real frequency
$\omega = \pi^{2}\nu\sigma_{1}/d^{2}$ and a horizontal wave number
$k = a/d$.  The horizontal phase velocity is
\begin{equation}
  v_{\mathrm{ph}} = \frac{\omega}{k}
  = \frac{\pi^{2}\nu\,\sigma_{1}}{a\,d} .
  \label{eq:rel-vph}
\end{equation}

\textbf{Requirement:} $v_{\mathrm{ph}} \le c$ for all physically
admissible overstable modes.

In the asymptotic regime $\Qrel\to\infty$, the oscillation frequency
scales as $|\sigma_{1}|\sim \vA^{(\mathrm{rel})}\,d/(\nu\pi^{2})$
(from~\eqref{eq:rel-gyration}) and $a\sim\Qrel^{1/6}$, so
\begin{equation}
  v_{\mathrm{ph}}
  \sim \frac{\vA^{(\mathrm{rel})}}{a_{c}\,d/d}
  = \frac{\vA^{(\mathrm{rel})}}{a_{c}} .
  \label{eq:rel-vph-asymp}
\end{equation}
Since $\vA^{(\mathrm{rel})} < c$ (cf.\ eq.~\eqref{eq:rel-Alfven}) and
$a_{c}\to\infty$ as $\Qrel\to\infty$, the phase velocity satisfies
$v_{\mathrm{ph}}\to 0$ in the strong-field limit.  Thus \emph{the
asymptotic overstable modes are automatically subluminal}.

For finite~$\Qrel$, the causality bound $v_{\mathrm{ph}} \le c$ imposes
the constraint
\begin{equation}
  \sigma_{1}^{2} \le \frac{c^{2}\,a^{2}}{\pi^{4}\nu^{2}} ,
  \label{eq:rel-causality-bound}
\end{equation}
which, combined with~\eqref{eq:rel-sigma1-sq}, gives a lower bound on
the wave number:
\begin{equation}
  a^{2} \ge \frac{\pi^{4}\nu^{2}}{c^{2}}\,
  \frac{1}{p_{1}^{(0)\,2}}\,
  \left[\frac{p_{2}^{(\mathrm{rel})}-p_{1}^{(\mathrm{rel})}}{1+p_{2}^{(\mathrm{rel})}}\,Q_{1}
        - (1+x)^{2}\right] .
  \label{eq:rel-a-causality-lower}
\end{equation}
Under terrestrial conditions ($\nu/c \sim 10^{-12}$\,s) this bound is
never restrictive, but in astrophysical settings (neutron star envelopes,
accretion discs near compact objects) where $\nu$ can be relativistically
large, the causality constraint may exclude long-wavelength overstable
modes.

The BDNK formalism guarantees that the
\emph{group velocity} of all dissipative modes is bounded by~$c$,
provided the transport coefficients and frame coefficients satisfy
the coupled BDNK inequalities
\citep{Bemfica2018,Kovtun2019,BemficaDisconziNoronhaKovtun2023}.
Unlike the Israel--Stewart approach, no separate group-velocity
bounds involving relaxation times need to be checked---the BDNK
strong hyperbolicity proof covers all modes simultaneously.
Since the characteristic equation has the same number of roots as
the classical one, there are no spurious relaxation modes whose
causality would need to be verified separately.
\end{causalitycheck}

%------------------------------------------------------------------
\subsection{Non-relativistic limit}
\label{subsec:rel-46-NR}

In the limit $c\to\infty$ (equivalently $h\to 1$, $\vA^{(\mathrm{rel})}
\to V_{A}$):
\begin{itemize}
  \item $\Qrel\to Q$, $\Rarel\to R$, recovering the classical
    definitions.
  \item $p_{1}^{(\mathrm{rel})}\to p_{1} = \kappa/\nu$, \;
    $p_{2}^{(\mathrm{rel})}\to p_{2} = \eta/\nu$.
  \item Equation~\eqref{eq:rel-char-full} $\to$ equation~(4-224).
  \item Equations~\eqref{eq:rel-sigma1-sq}--\eqref{eq:rel-R1-overstab}
    $\to$ equations~(4-230)--(4-231).
  \item The BDNK causality bounds become trivially satisfied.
  \item The BDNK frame coefficients vanish at $\order{v^{2}/c^{2}}$.
\end{itemize}
All classical results of \S\S\,45--46 are thus recovered.


%% =============================================================
%% Astrophysical application: pi^2 Q law verification for NS
%% =============================================================
\subsection{Astrophysical application: the $\pi^{2}Q$ law in neutron
  star magneto-convection}
\label{subsec:rel-46-NS-application}

The $\pi^{2}Q$ law~\eqref{eq:rel-pi2Q-law} has a direct consequence
for thermal convection in neutron star interiors.  In the strong-field
regime characteristic of magnetars ($\Qrel \gg 1$), the critical
temperature gradient for the onset of convection is determined
entirely by the magnetic field strength and the electrical conductivity,
independent of viscosity (equation~\eqref{eq:rel-beta-c-inviscid}).

For a magnetar core with $B \sim 10^{15}$--$10^{16}$~G,
$\sigma_{\mathrm{e}} \sim 10^{24}$~s$^{-1}$
\citep{YakovlevPethick2004},
$\rho \sim 10^{14}$~g\,cm$^{-3}$, and layer depth $d \sim 10^{5}$~cm,
the classical Chandrasekhar number is $Q \sim 10^{8}$--$10^{12}$.
The relativistic correction reduces this by a factor
$\rho c^{2}/\enthalpy \approx 0.85$--$0.95$ for nuclear-density matter.

Figure~\ref{fig:pi2Q-verification} demonstrates the convergence of
$\Rarel/\Qrel \to \pi^{2}$ in the strong-field limit for both free
and rigid boundary conditions.

\begin{figure}[htbp]
  \centering
  \includegraphics[width=0.95\textwidth]{fig_pi2Q_law_verification.pdf}
  \caption{Numerical verification of the $\pi^{2}Q$ law.
    \emph{Left:} Critical Rayleigh number $\Rarel$ versus
    Chandrasekhar number $\Qrel$ for free boundaries (solid blue),
    rigid boundaries (dashed purple), and relativistic free
    boundaries with $\vA^{2}/c^{2} = 0.05$ (dash-dotted red).
    \emph{Right:} The ratio $\Rarel/\Qrel$ converges to $\pi^{2}$
    (dotted grey line) for all boundary conditions as $\Qrel\to\infty$,
    confirming the universality of the $\pi^{2}Q$ law in the
    relativistic theory.}
  \label{fig:pi2Q-verification}
\end{figure}

The invariant~\eqref{eq:rel-invariant-result} provides an additional
observational diagnostic: the balance between magnetic and kinetic
energy perturbations at marginal stability is shifted by the factor
$h = \enthalpy/(\rho c^{2})$ relative to the classical prediction.
For a radiation-dominated neutron star envelope where
$\enthalpy \approx 4p$, the invariant can be a factor of
$\sim 1.3$--$4$ larger than the classical value, depending on the
equation of state \citep{OertelHempel2017}.
